% =====================================================================
%  THÈME 3 — Arrondir & opérations — EXERCICES — Niveau FACILE
% =====================================================================
\documentclass[11pt,a4paper]{article}
\usepackage[utf8]{inputenc}
\usepackage[T1]{fontenc}
\usepackage{lmodern}
\usepackage[french]{babel}
\usepackage{amsmath,amssymb,mathtools}
\usepackage{icomma}
\usepackage{siunitx}
\sisetup{output-decimal-marker={,},per-mode=symbol}
\usepackage[version=4]{mhchem}
\usepackage{xcolor}
\usepackage[most]{tcolorbox}
\usepackage{enumitem}
\usepackage{fancyhdr}
\usepackage[a4paper,margin=2.2cm,headheight=26pt,headsep=8pt]{geometry}

\definecolor{primary}{RGB}{150,98,162}
\definecolor{primarydark}{RGB}{92,52,110}
\newcommand{\cs}{chiffre significatif}
\newcommand{\CS}{chiffres significatifs}

\pagestyle{fancy}\fancyhf{}
\fancyhead[L]{\small\textcolor{primarydark}{\textbf{Fiche 1} — Chiffres significatifs}}
\fancyhead[R]{\small\textcolor{primarydark}{\textsc{Exercices · Facile · Thème 3}}}
\fancyfoot[C]{\small\thepage}
\renewcommand{\headrulewidth}{0.8pt}
\renewcommand{\headrule}{\hbox to\headwidth{\color{primary}\leaders\hrule height \headrulewidth\hfill}}

\newtcolorbox[auto counter]{exo}[1][]{enhanced,breakable,boxrule=0.8pt,arc=2pt,
  colback=primary!4,colframe=primary,left=3mm,right=3mm,top=2mm,bottom=2mm,
  fonttitle=\bfseries,coltitle=white,
  title={Exercice~\thetcbcounter\ifx\relax#1\relax\relax\else\ ---\ \textmd{\itshape #1}\fi}}

\setlength{\parindent}{0pt}\setlength{\parskip}{3pt}

\begin{document}

\begin{center}
{\Large\bfseries\color{primarydark} Exercices — Niveau Facile}\\[1pt]
{\color{primary} Thème 3 : arrondir un résultat}
\end{center}
\vspace{1mm}
\textit{On travaille uniquement le geste d'arrondi (les opérations viennent au niveau suivant).}
\vspace{2mm}

\begin{exo}[Arrondir à un nombre de \CS\ donné]
Arrondis chaque nombre au nombre de \CS\ demandé.
\begin{enumerate}[label=\alph*),leftmargin=1.6em,itemsep=1pt]
  \item $\num{3,148}$ à $3$~CS \qquad b) $\num{3,142}$ à $3$~CS \qquad c) $\num{0,06237}$ à $2$~CS
  \item[] d) $\num{124,7}$ à $3$~CS \qquad e) $\num{7,8451}$ à $2$~CS
\end{enumerate}
\end{exo}

\begin{exo}[Le premier chiffre supprimé]
Pour chaque arrondi demandé, indique d'abord le \emph{premier chiffre supprimé}, dis si l'on arrondit
par défaut ou par excès, puis donne le résultat.
\begin{enumerate}[label=\alph*),leftmargin=1.6em,itemsep=1pt]
  \item $\num{0,4271}$ à $2$~CS \qquad b) $\num{58,96}$ à $3$~CS \qquad c) $\num{1,2350}$ à $3$~CS
\end{enumerate}
\end{exo}

\begin{exo}[Attention aux retenues en cascade]
Arrondis ; une retenue peut se propager. Pense à \emph{conserver} les zéros nécessaires au bon nombre
de \CS.
\begin{enumerate}[label=\alph*),leftmargin=1.6em,itemsep=1pt]
  \item $\num{2,996}$ à $3$~CS \qquad b) $\num{9,98}$ à $2$~CS \qquad c) $\num{0,7996}$ à $3$~CS
\end{enumerate}
\end{exo}

\begin{exo}[Arrondir à un nombre de \emph{décimales} donné]
Arrondis au nombre de décimales indiqué (chiffres après la virgule).
\begin{enumerate}[label=\alph*),leftmargin=1.6em,itemsep=1pt]
  \item $\num{255,507}$ à $1$ décimale \qquad b) $\num{0,1149}$ à $2$ décimales \qquad c) $\num{18,016}$ à $2$ décimales
\end{enumerate}
\end{exo}

\begin{exo}[Arrondir de grands nombres]
Arrondis au nombre de \CS\ demandé, puis écris le résultat en notation scientifique.
\begin{enumerate}[label=\alph*),leftmargin=1.6em,itemsep=1pt]
  \item $\num{45678}$ à $2$~CS \qquad b) $\num{30960}$ à $3$~CS \qquad c) $\num{1249}$ à $2$~CS
\end{enumerate}
\end{exo}

\end{document}
